Mechanistic models are central to quantitative understanding and optimisation of Chinese hamster ovary (CHO) cell culture processes, but their utility is often restricted by parameter sets calibrated for specific cell lines, scales, or operating conditions. In this study, we present the application of the Ensemble Kalman Filter (EnKF) to bioprocessing, introducing an  ensemble-based framework for dual state and parameter estimation that enables mechanistic model adaptation across distinct systems.\textcolor{blue}{The EnKF recursively assimilates process measurements to update uncertain kinetic parameters and predict system states. This allows a model calibrated for one system to be transferred to another without reparametrisation, using only a single experimental dataset.} The evolving parameter ensembles provide a time-resolved sensitivity analysis that identifies which parameters have dominant influence under new process conditions and when their effects become significant. \textcolor{blue}{The framework was evaluated across six CHO cell datasets spanning different scales, cell lines, temperatures, and feeding strategies. It accurately reconstructed system dynamics and showed progressive improvement in long-term predictions as data accumulated.} By maintaining full mechanistic transparency while flexibly adapting to new data, the EnKF offers a practical route for knowledge transfer across systems, strengthening the role of mechanistic modelling in data-informed bioprocess understanding and control.


